\chapter{Reading a Macro Calendar and a Rates Screen}\label{ch:m2:reading-a-macro-calendar-and-a-rates-screen}

On most Fridays at the start of a month, at 08:30 in New York, the Bureau of Labor
Statistics publishes the employment report. From January 2023 to September 2026 the
2-year Treasury yield moved, on the 44 days it came out, by 10.2 basis points on average,
up or down; on the other 887 trading days, by 4.6. The rates market lives by a calendar of
such releases and central-bank meetings, and a trader reads it the way a pilot reads a
weather map: where the turbulence will be, how bad it can get, and how exposed the
position is when it arrives. This chapter closes the book on that craft: the calendar,
consensus and surprise, positioning into events, the shapes a curve can take and how to
trade them, and what a rates screen shows at a glance.

\section{The calendar}

\begin{definition}[Economic release]\label{def:m2:reading-a-macro-calendar-and-a-rates-screen:release}
An \emph{economic release}\index{economic release} is the scheduled publication of an
official statistic, such as employment or consumer prices, at a date and time announced in
advance by the agency that produces it.
\end{definition}

In the United States, among the releases that move rates, the employment report and the
consumer price index both come at 08:30 Eastern time; the Federal Reserve's policy
decisions come at the end of its eight scheduled meetings a year. The calendar for
September 2026 shows how they interlock
(\cref{fig:m2:reading-a-macro-calendar-and-a-rates-screen:calendar}): the August
employment report on Friday 4 September, the August CPI on Friday 11 September, and the
Federal Open Market Committee on 15 and 16 September, one of the four meetings of the year
with the Committee's economic projections.

\begin{omfigure}
\centering
\begin{tikzpicture}[font=\small,x=0.62cm]
\fill[omThm!12] (5,-0.2) rectangle (18,0.2);
\draw[thick,-{Stealth[length=2.5mm]}] (0.5,0) -- (21.5,0);
\foreach \d in {1,2,3,4,5,6,7,8,9,10,11,12,13,14,15,16,17,18,19,20,21} \draw (\d,-0.08) -- (\d,0.08);
\foreach \d in {1,4,5,7,11,14,15,16,17,21} \node[font=\scriptsize,above] at (\d,0.2) {\d};
\node[font=\footnotesize,text=omThm] at (11.5,1.0) {FOMC blackout: Saturday 5 to Thursday 17 September};
\node[draw,thick,fill=omDef!10,rounded corners=2pt,align=center,font=\footnotesize] (p) at (4,-1.9) {Fri 4, 08:30\\employment report};
\node[draw,thick,fill=omProp!10,rounded corners=2pt,align=center,font=\footnotesize] (c) at (11,-1.9) {Fri 11, 08:30\\CPI};
\node[draw,thick,fill=omThm!10,rounded corners=2pt,align=center,font=\footnotesize] (f) at (16.5,-1.9) {Tue 15 -- Wed 16\\FOMC, with projections};
\draw[-{Stealth[length=2mm]}] (p.north) -- (4,-0.1);
\draw[-{Stealth[length=2mm]}] (c.north) -- (11,-0.1);
\draw[-{Stealth[length=2mm]}] (f.north) -- (16,-0.1);
\end{tikzpicture}

\omcaption{The US rates calendar for the first three weeks of September 2026: the
employment report and CPI at 08:30 Eastern time, and the FOMC meeting of 15 and 16
September with its \omterm{def:m2:reading-a-macro-calendar-and-a-rates-screen:blackout}{blackout period}, from the second Saturday before the meeting to the day
after it. Dates from the BLS and Federal Reserve calendars.}\label{fig:m2:reading-a-macro-calendar-and-a-rates-screen:calendar}
\end{omfigure}

\begin{definition}[Blackout period]\label{def:m2:reading-a-macro-calendar-and-a-rates-screen:blackout}
A \emph{blackout period}\index{blackout period} is the interval around a policy meeting
during which a central bank's policymakers refrain from public comment on the economy and
monetary policy.
\end{definition}

The FOMC's rule, reaffirmed in January 2026, starts the blackout at midnight on the second
Saturday before a meeting and ends it at 23:59 on the day after: for a Tuesday--Wednesday
meeting, from the Saturday ten days earlier to the Thursday. In September 2026 the CPI
fell inside it, so the market had to price the report without hearing from the
Committee's members. Calendars also carry holidays, auctions (chapter 4) and index
rebalancings (chapter 17); a trader's calendar is the union.

\begin{dated}{2026-09}{US rates calendar}\label{dat:m2:reading-a-macro-calendar-and-a-rates-screen:calendar}
Employment report 2026: 08:30 ET, e.g. 4 September, 2 October, 6 November, 4 December; no
report was published in October 2025, during the lapse in federal appropriations. CPI 2026:
08:30 ET, e.g. 11 September, 14 October. FOMC 2026: 27--28 January, 17--18 March, 28--29
April, 16--17 June, 28--29 July, 15--16 September, 27--28 October, 8--9 December (March,
June, September and December with projections). FOMC blackout rule reaffirmed 27 January 2026.
\end{dated}

\section{Consensus and surprise}

\begin{definition}[Consensus forecast, data surprise]\label{def:m2:reading-a-macro-calendar-and-a-rates-screen:surprise}
The \emph{consensus forecast}\index{consensus forecast} of a release is the median of
economists' forecasts collected before it. The \emph{data surprise}\index{data surprise}
is the released value minus the consensus, often divided by the standard deviation of
past surprises to make releases comparable.
\end{definition}

Markets price the consensus before the release; what moves them is the surprise. The
first measure of a release's importance is how much markets move on its days
(\cref{fig:m2:reading-a-macro-calendar-and-a-rates-screen:moves}). On payroll days from
2023 to 2026 the 2-year yield moved 2.2 times as much as on other days, the 10-year 1.7
times, and the 2s10s slope 1.4 times: the report is about the path of policy, which the
front end prices most directly.

\begin{omfigure}
\begin{tikzpicture}
\begin{axis}[width=11cm,height=4.8cm,ybar,bar width=18pt,
  ylabel={mean absolute daily change (bp)},
  tick label style={font=\small},label style={font=\small},
  xtick=data,xticklabels={2-year,10-year,2s10s},enlarge x limits=0.3,
  ymin=0,ymax=12,grid=major,grid style={gray!25},
  legend columns=2,legend style={font=\footnotesize,at={(0.5,-0.22)},anchor=north,draw=none,fill=none,/tikz/every even column/.append style={column sep=8pt}},
  nodes near coords,nodes near coords style={font=\scriptsize},
  every node near coord/.append style={/pgf/number format/fixed,/pgf/number format/fixed zerofill,/pgf/number format/precision=1},
  table/col sep=comma]
\addplot[fill=omDef!60,draw=omDef] table[x=k,y=payroll]{figdata/markets-2/31-reading-a-macro-calendar-and-a-rates-screen/eventmoves.csv};
\addplot[fill=gray!35,draw=gray] table[x=k,y=other]{figdata/markets-2/31-reading-a-macro-calendar-and-a-rates-screen/eventmoves.csv};
\legend{44 payroll days,887 other days}
\end{axis}
\end{tikzpicture}

\omcaption{Mean absolute daily change in the 2- and 10-year Treasury yields and in the
2s10s slope on employment-report days and on other days, January 2023 to September 2026.
Data: Federal Reserve H.15 via FRED (DGS2, DGS10); release dates from the BLS
archive.}\label{fig:m2:reading-a-macro-calendar-and-a-rates-screen:moves}
\end{omfigure}

The second measure is the sensitivity: how many basis points a yield moves per unit of
standardised surprise, estimated by regressing moves on surprises over many releases.
Consensus data are sold by vendors, so the tutorial shows the method on synthetic data:
with 200 releases whose true sensitivities are 8 and 5 basis points for the 2- and 10-year
yields, the regression finds 8.4 and 5.1, with standard errors of 0.2.

\section{Positioning into events}

A position that is harmless on a quiet day can be dangerous over a release. The question
before an event is the size of the move one could face, which the history of event days
answers, and whether one wants the exposure at all: a trader with a view expresses it; a
market maker reduces risk or widens quotes into the number; a portfolio that is not meant
to bet on data trims its duration or hedges the front end.

\begin{definition}[Steepener, flattener, curve fly]\label{def:m2:reading-a-macro-calendar-and-a-rates-screen:steep}
A \emph{steepener}\index{steepener} is a position that gains when the yield curve
steepens, typically long a shorter bond and short a longer one with equal DV01; a
\emph{flattener}\index{flattener} is the opposite. A \emph{curve fly}\index{curve fly} is
a position in three maturities, long or short the middle against the two wings, sized to be
neutral to the level and slope of the curve; its price is measured by twice the middle
yield minus the two wings.
\end{definition}

A DV01-neutral \omterm{def:m2:reading-a-macro-calendar-and-a-rates-screen:steep}{steepener} is insensitive to parallel moves: it earns or loses the change in
the slope times its DV01 per leg. With the chapter 3 bond calculator, an illustrative
10-year note with a coupon of 4.875\% at a yield of 4.96\% has a DV01 of 0.0772 per 100, a
2-year note of 4.625\% at 4.71\% one of 0.0183. Short USD~100 million of the 10-year has a
DV01 of USD~77\,168 per basis point; the 2-year leg that matches it is USD~421 million.

\section{Reading curve shapes}

\begin{definition}[Bull steepening, bear flattening]\label{def:m2:reading-a-macro-calendar-and-a-rates-screen:bull}
A move of the yield curve is \emph{bull steepening}\index{bull steepening} when yields fall
and the curve steepens, the short end falling more, and \emph{bear flattening}\index{bear flattening}
when yields rise and the curve flattens, the short end rising more. Bull flattening and
bear steepening are the other two combinations.
\end{definition}

Most payroll days move the curve in one of the two ways defined above, the front end
leading (\cref{fig:m2:reading-a-macro-calendar-and-a-rates-screen:scatter}): of the 44 days,
14 were \omterm{def:m2:reading-a-macro-calendar-and-a-rates-screen:bull}{bull steepening} and 20 \omterm{def:m2:reading-a-macro-calendar-and-a-rates-screen:bull}{bear flattening}. On 2 August 2024 the 2-year fell 28 basis
points and the 10-year 19; on 4 October 2024 the 2-year rose 23 and the 10-year 13. A weak
report pulls expected policy rates down, which lowers the front end most; a strong one
does the reverse. Moves of the other two kinds, in which the long end leads, often point
to something other than the policy path: term premium, supply, inflation expectations.

\begin{omfigure}
\begin{tikzpicture}
\begin{axis}[width=9cm,height=7.2cm,
  xlabel={change in 2-year yield (bp)},ylabel={change in 10-year yield (bp)},
  tick label style={font=\small},label style={font=\small},
  xmin=-32,xmax=32,ymin=-26,ymax=26,grid=major,grid style={gray!25},
  legend columns=2,legend style={font=\footnotesize,at={(0.5,-0.2)},anchor=north,draw=none,fill=none,/tikz/every even column/.append style={column sep=8pt}},
  scatter/classes={bull-steepening={mark=*,omDef},bear-flattening={mark=square*,omThm},bull-flattening={mark=triangle*,omProp},bear-steepening={mark=diamond*,gray}},
  table/col sep=comma]
\addplot[scatter,only marks,mark size=2pt,scatter src=explicit symbolic] table[x=d2,y=d10,meta=kind]{figdata/markets-2/31-reading-a-macro-calendar-and-a-rates-screen/payrollmoves.csv};
\addplot[black,thin,dashed,domain=-26:26,forget plot] {x};
\addplot[gray,thin,forget plot] coordinates {(-32,0) (32,0)};
\addplot[gray,thin,forget plot] coordinates {(0,-26) (0,26)};
\legend{bull steepening,bear flattening,bull flattening,bear steepening}
\end{axis}
\end{tikzpicture}

\omcaption{Changes in the 2- and 10-year Treasury yields on the 44 employment-report days of
January 2023 to September 2026, classified by the average of the two changes (bull or
bear) and the change in 2s10s (steepening above the dashed diagonal, flattening below it).
\omterm{def:m2:reading-a-macro-calendar-and-a-rates-screen:bull}{Bull steepening}, on the left above the diagonal, and \omterm{def:m2:reading-a-macro-calendar-and-a-rates-screen:bull}{bear flattening}, on the right below
it, are 34 of the 44. Data: H.15 via FRED; BLS release dates.}\label{fig:m2:reading-a-macro-calendar-and-a-rates-screen:scatter}
\end{omfigure}

\section{A rates screen at a glance}

A rates trader's screen condenses the market into a few numbers: the level of the curve;
its slopes (2s10s, 5s30s); its curvature, the 2s5s10s fly; the \omterm{def:m2:short-term-interest-rate-futures:path}{implied policy path} of
chapter 8; \omterm{def:m2:interest-rate-swaps:spread}{swap spreads} and the \omterm{def:m2:fx-swaps-forwards-and-the-cross-currency-basis:basis}{cross-currency basis} of chapters 9 and 16; the volatility
of chapter 13; and the calendar. Read together they say where the market is, what it
expects, and what has changed. From September 2025 to September 2026 the Treasury curve
rose and flattened: the 2-year rose from 3.61\% to 4.71\%, the 10-year from 4.15\% to 4.96\%,
and 2s10s fell from 54 to 25 basis points, while the 2s5s10s fly moved from $-34$ to $-1$
(\cref{fig:m2:reading-a-macro-calendar-and-a-rates-screen:curves}).

\begin{omfigure}
\begin{tikzpicture}
\begin{axis}[width=11cm,height=5cm,
  xlabel={maturity (years)},ylabel={yield (\%)},
  tick label style={font=\small},label style={font=\small},
  xmode=log,log basis x=10,xtick={2,5,10,30},xticklabels={2,5,10,30},
  xmin=1.7,xmax=35,ymin=3.4,ymax=5.5,grid=major,grid style={gray!25},
  legend columns=3,legend style={font=\footnotesize,at={(0.5,-0.3)},anchor=north,draw=none,fill=none,/tikz/every even column/.append style={column sep=8pt}},
  table/col sep=comma]
\addplot[gray,very thick,mark=*] table[x=tenor,y=d20250922]{figdata/markets-2/31-reading-a-macro-calendar-and-a-rates-screen/curves.csv};
\addplot[omProp,very thick,mark=*] table[x=tenor,y=d20260323]{figdata/markets-2/31-reading-a-macro-calendar-and-a-rates-screen/curves.csv};
\addplot[omDef,very thick,mark=*] table[x=tenor,y=d20260922]{figdata/markets-2/31-reading-a-macro-calendar-and-a-rates-screen/curves.csv};
\legend{22 Sep 2025,23 Mar 2026,22 Sep 2026}
\end{axis}
\end{tikzpicture}

\omcaption{The Treasury curve at 2, 5, 10 and 30 years on three dates six months apart:
over the year it rose and flattened, a \omterm{def:m2:reading-a-macro-calendar-and-a-rates-screen:bull}{bear flattening} led by the 2-year. Data: Federal
Reserve H.15 via FRED (DGS2, DGS5, DGS10, DGS30).}\label{fig:m2:reading-a-macro-calendar-and-a-rates-screen:curves}
\end{omfigure}

\begin{dated}{2026-09}{A Treasury screen}\label{dat:m2:reading-a-macro-calendar-and-a-rates-screen:screen}
Constant-maturity yields on 22 September 2026: 2-year 4.71\%, 5-year 4.83\%, 10-year 4.96\%,
30-year 5.29\%; 2s10s 25 basis points, 5s30s 46, 2s5s10s fly $-1$. A year earlier: 3.61, 3.71,
4.15 and 4.77\%; 54, 106 and $-34$.
\end{dated}

\section{Tutorial: a release calendar and a curve screen}\label{tut:m2:reading-a-macro-calendar-and-a-rates-screen:macrocal}

\begin{tutorial}
\textbf{Goal.} Build a release calendar with the blackout, measure moves on event days,
estimate surprise sensitivities, summarise the curve and size a DV01-neutral trade.
\textbf{End state:} \cref{fig:m2:reading-a-macro-calendar-and-a-rates-screen:calendar,fig:m2:reading-a-macro-calendar-and-a-rates-screen:moves,fig:m2:reading-a-macro-calendar-and-a-rates-screen:scatter,fig:m2:reading-a-macro-calendar-and-a-rates-screen:curves}
and the numbers of the weekend problem.
\begin{tutsteps}
\item \textbf{The calendar}: events, the blackout rule, and flags.
\omcode{code/firm/macrocal/firm_macrocal.py}{13}{29}{Events, the FOMC blackout and the calendar.}\label{lst:m2:reading-a-macro-calendar-and-a-rates-screen:cal}
\item \textbf{The screen}: slopes and fly, move classification, event-day moves,
DV01-neutral sizing and P\&L.
\omcode{code/firm/macrocal/firm_macrocal.py}{32}{61}{Curve summary, classification, event moves and sizing.}\label{lst:m2:reading-a-macro-calendar-and-a-rates-screen:screen}
\item \textbf{Surprises}: the sensitivity of yields to standardised surprises.
\omcode{code/firm/macrocal/firm_macrocal.py}{64}{72}{Sensitivity to surprises by least squares.}\label{lst:m2:reading-a-macro-calendar-and-a-rates-screen:sens}
\item \textbf{Run} \texttt{macro\_demo.event\_stats()}, \texttt{macro\_demo.problem()},
\texttt{macro\_demo.surprise\_regression()} and \texttt{fig\_macro.py}.
\end{tutsteps}
\textbf{What to change next.} Add the CPI release dates and compare CPI days with
payroll days; then fit the 2-year's sensitivity to payroll surprises if you have
consensus data, and compare it with the average move.
\end{tutorial}

\section{Build: event calendar and curve summary}\label{bld:m2:reading-a-macro-calendar-and-a-rates-screen:macrocal}

\begin{build}
\bfield{Purpose} The miniature firm's morning note and its risk limits around events: what
is released when, which meetings are in blackout, how the curve looks, and how large the
curve positions are in DV01.
\bfield{Interface} \texttt{Event(day, time, name)}; \texttt{blackout(meeting\_start,
meeting\_end)}; \texttt{calendar(events, meetings)}; \texttt{curve\_summary(y2, y5, y10,
y30)}; \texttt{classify(d2, d10)}; \texttt{event\_day\_moves(dates, y, events)};
\texttt{dv01\_neutral(dv01\_short\_leg, dv01\_long\_leg, notional\_short)};
\texttt{steepener\_pnl(dv01, d2\_bp, d10\_bp)}; \texttt{sensitivity(surprises, moves)}.
\bfield{Rules} Yields in percent, changes in basis points; blackout from the second Saturday
before the meeting to the day after (the federal-holiday exception not implemented);
first-order P\&L for curve trades; OLS with an intercept for sensitivities.
\bfield{Acceptance tests} \texttt{code/firm/macrocal/tests/}: the September 2026 and January
2026 blackouts; the calendar flags CPI inside the blackout and payrolls outside; slopes and
fly of 22 September 2026; the four classifications; event-day means, DV01-neutral notional
and \omterm{def:m2:reading-a-macro-calendar-and-a-rates-screen:steep}{steepener} P\&L on small cases; the regression recovers a known slope.
\bfield{Stretch} Consensus and surprise series from a data vendor; intraday moves around
releases; the \omterm{def:m2:short-term-interest-rate-futures:path}{implied policy path} before and after each meeting (chapter 8); a morning
screen that combines all the chapters' builds.
\end{build}

\begin{omsources}
\item Bureau of Labor Statistics, release schedules for the Employment Situation and the
CPI; archive of Employment Situation releases.
\item Federal Reserve, FOMC calendars; FOMC Policy on External Communications of Committee
Participants (reaffirmed January 2026).
\item Board of Governors of the Federal Reserve System, H.15 selected interest rates, via
FRED (DGS2, DGS5, DGS10, DGS30).
\end{omsources}

\section{Exercises}

\begin{exercise}[$\star$]\label{exo:m2:reading-a-macro-calendar-and-a-rates-screen:1}
The FOMC meets on Tuesday 27 and Wednesday 28 October 2026. When does its blackout begin
and end?
\end{exercise}

\begin{exercise}[$\star$]\label{exo:m2:reading-a-macro-calendar-and-a-rates-screen:2}
The 2-year yield falls 10 basis points and the 10-year 4. Name the move.
\end{exercise}

\begin{exercise}[$\star$]\label{exo:m2:reading-a-macro-calendar-and-a-rates-screen:3}
Compute 2s10s, 5s30s and the 2s5s10s fly on 22 September 2025.
\end{exercise}

\begin{exercise}[$\star\star$]\label{exo:m2:reading-a-macro-calendar-and-a-rates-screen:4}
Why does the employment report move the 2-year yield more than the 10-year?
\end{exercise}

\begin{exercise}[$\star\star$]\label{exo:m2:reading-a-macro-calendar-and-a-rates-screen:5}
A DV01-neutral \omterm{def:m2:reading-a-macro-calendar-and-a-rates-screen:steep}{steepener} has a DV01 of USD~50\,000 per leg. What does it earn if the 2-year
falls 15 basis points and the 10-year 5? If both fall 10?
\end{exercise}

\begin{exercise}[$\star\star$]\label{exo:m2:reading-a-macro-calendar-and-a-rates-screen:6}
Why is a pass through a release, with no position, sometimes the best position?
\end{exercise}

\begin{exercise}[$\star\star\star$]\label{exo:m2:reading-a-macro-calendar-and-a-rates-screen:7}
\emph{Coding.} With \texttt{sensitivity} and the synthetic surprises, how many releases
are needed for the standard error of the 2-year sensitivity to fall to 0.5 basis points?
Check the answer by simulation.
\end{exercise}

\begin{exercise}[$\star\star\star$]\label{exo:m2:reading-a-macro-calendar-and-a-rates-screen:8}
\emph{Find the flaw.} ``The 2-year moves 10 basis points on payroll days; my position has a
DV01 of USD~100\,000; so my payroll risk is USD~1 million.'' Correct it.
\end{exercise}

\section{Problem: Payrolls Friday}

\begin{problem}[{Weekend problem --- a steepener into the employment report}]\label{pb:m2:reading-a-macro-calendar-and-a-rates-screen:1}
On a Thursday evening a trader expects a weak employment report and wants a curve position
that gains from it without betting on the level of rates. It sells USD~100 million of an
illustrative 10-year note (coupon 4.875\%, yield 4.96\%) and buys a 2-year note (coupon
4.625\%, yield 4.71\%), DV01-neutral, settling on 22 September 2026.

\medskip\noindent\textbf{Part I --- Sizing.}
\begin{enumerate}
\item What are the two notes' DV01s per 100 of face?
\item What is the DV01 of the 10-year leg?
\item How much of the 2-year does the trader buy?
\item Why DV01-neutral rather than equal notional?
\item What does the position gain or lose if all yields move by the same amount?
\end{enumerate}

\medskip\noindent\textbf{Part II --- The report.}
\begin{enumerate}[resume]
\item Suppose the day repeats 2 August 2024: the 2-year falls 28 basis points and the 10-year
19. What is the P\&L to first order? Fully repriced?
\item Suppose instead 4 October 2024: the 2-year rises 23 and the 10-year 13.
\item Name each of the two moves.
\item How often, on the 44 payroll days of 2023--2026, did the curve move one of these two
ways?
\item What does the difference between the first-order and repriced P\&L come from?
\end{enumerate}

\medskip\noindent\textbf{Part III --- The calendar.}
\begin{enumerate}[resume]
\item The report comes on 4 September 2026. Is the FOMC in blackout?
\item The CPI comes a week later. Why does the blackout matter for how the market takes it?
\item How much more does the 2-year move on payroll days than on other days?
\item What else on the calendar could move the position that week?
\item How would you size the position to lose at most USD~500\,000 on a typical bad payroll
day?
\end{enumerate}

\medskip\noindent\textbf{Part IV --- Judgement.}
\begin{enumerate}[resume]
\item Why might the \omterm{def:m2:reading-a-macro-calendar-and-a-rates-screen:steep}{steepener} lose even if the report is weak?
\item What would you watch on the screen before and after the release?
\item How does the trade relate to the \omterm{def:m2:short-term-interest-rate-futures:path}{implied policy path} of chapter 8?
\item State the \emph{named result}: the DV01-neutral 2s10s position and its P\&L under
each of the two moves.
\item In one sentence: what does a DV01-neutral \omterm{def:m2:reading-a-macro-calendar-and-a-rates-screen:steep}{steepener} bet on?
\end{enumerate}
\end{problem}

\section{Interview questions}

\begin{interviewq}[$\star$ \iqroles{trader}]\label{iq:m2:reading-a-macro-calendar-and-a-rates-screen:1}
What are the most important US data releases for rates, and when do they come out?
\end{interviewq}

\begin{interviewq}[$\star$ \iqroles{trader, researcher}]\label{iq:m2:reading-a-macro-calendar-and-a-rates-screen:2}
Define \omterm{def:m2:reading-a-macro-calendar-and-a-rates-screen:bull}{bull steepening} and \omterm{def:m2:reading-a-macro-calendar-and-a-rates-screen:bull}{bear flattening}, and give a cause of each.
\end{interviewq}

\begin{interviewq}[$\star\star$ \iqroles{trader}]\label{iq:m2:reading-a-macro-calendar-and-a-rates-screen:3}
How do you construct a 2s10s \omterm{def:m2:reading-a-macro-calendar-and-a-rates-screen:steep}{steepener}, and what are its risks?
\end{interviewq}

\begin{interviewq}[$\star\star$ \iqroles{researcher}]\label{iq:m2:reading-a-macro-calendar-and-a-rates-screen:4}
How would you measure the market's sensitivity to a data release?
\end{interviewq}

\begin{interviewq}[$\star\star$ \iqroles{trader, risk}]\label{iq:m2:reading-a-macro-calendar-and-a-rates-screen:5}
What would you put on a one-page rates screen, and why?
\end{interviewq}

\begin{interviewq}[$\star\star\star$ \iqroles{developer}]\label{iq:m2:reading-a-macro-calendar-and-a-rates-screen:6}
Design the system that ingests a release at 08:30:00 and updates every desk's risk and
quotes within a second.
\end{interviewq}
